\section{Responsibility Without Artificial Persons}
\label{sec:definitions}

The cascade maps control exercised by people and organizations over a deployed system. That method requires one stable premise: operational performance does not silently change the legal cast of characters. A system may generate a promise, recommendation, denial, image, or instruction. Whether the effect is attributable to a company, professional, public body, user, or another recognized person depends on the governing law and the deployment facts. The system's causal contribution does not, by itself, make it the promisor, tortfeasor, author, fiduciary, or constitutional claimant.

This is a status rule, not a universal definition of AI and not a judgment about the technology's worth. Regulatory triggers should vary with sector and risk. Functional classifications should vary with doctrine. The invariant is narrower: unless positive law creates and equips a new juridical status, the system does not independently hold the claims, duties, powers, liabilities, assets, procedural capacity, and exposure to remedies necessary for adjudication. That premise directs courts back to the deployment rather than toward an artificial liability sink.

\subsection{Legal Status Is an Institutional Structure}

Legal personality begins with positions, not traits. Hohfeld's analysis remains useful because it disaggregates the word ``right'' into relations with distinct correlatives and opposites.\lrfootnote{See \artcite[30--58]{hohfeld1913}. Hohfeldian ``liability'' means exposure to another's legal power, not necessarily tort liability or capacity for punishment. The table is diagnostic; it does not imply that every legal person personally exercises every incident. Its remedy row is this Article's enforcement-oriented extension of the four Hohfeldian pairs.} Adapted to an AI dispute, the structure exposes what generated performance leaves unresolved:

\begin{table}[htbp]
\centering
\caption{Normative Positions and Artificial Performance}
\label{tab:hohfeld}
\small
\begin{tabularx}{\textwidth}{>{\raggedright\arraybackslash}p{0.15\textwidth} >{\raggedright\arraybackslash}p{0.16\textwidth} >{\raggedright\arraybackslash}p{0.25\textwidth} >{\raggedright\arraybackslash}X}
\toprule
Position & Correlative & Legal operation & What performance alone cannot establish \\
\midrule
Claim & Duty & One subject can demand conduct from another. & That the system owns the claim or can release, assign, and enforce it. \\
Privilege & No-right & One subject may act free of another's claim to the contrary. & That the system is the beneficiary of the legal permission. \\
Power & Liability & One subject can alter a legal relation by a recognized act. & That generated assent, notice, waiver, or transfer is the system's juridical act. \\
Immunity & Disability & One subject is protected from another's attempted alteration. & That the system holds a constitutional, statutory, or private-law immunity. \\
Remedy & Exposure & Institutions can order relief against assets, conduct, or status. & What property answers, who appears, and how judgment can be enforced or sanction experienced. \\
\bottomrule
\end{tabularx}
\end{table}

A model can produce factual effects at every row. It can classify an applicant, send a notice, trigger a payment, recommend a denial, or generate assent-like text. The legal effect exists, however, because law identifies a subject and supplies an attribution rule. The bank's automated process transfers funds under authority attributable to the bank or customer. The agency's automated notice matters because the agency used it. The denial is reviewable because a public body or employer made the system part of its decision process.

Competence therefore cannot be the legal test. Law recognizes newborn children and people with profound cognitive disabilities without conditioning human status on linguistic or planning performance. It also constitutes juridical persons that possess no consciousness. Corporations own property, sue and are sued, contract, commit torts and crimes, and survive changes in membership because statutes and common law identify their organs, attribute acts and knowledge, partition assets, authorize service, and make judgments enforceable.\lrfootnote{See \casecite[636]{dartmouth1819}; \casecite[492--95]{newYorkCentral1909}; \artcite{deweyCorporatePersonality1926}; \bookcite{kurkiLegalPersonhood2019}.} Corporate personality is an institutional achievement, not an intelligence benchmark for machines.

The point also defeats a recurring circularity. If a chatbot is said to possess a legal power because it generated promise-like words, its capacity is inferred from the very act whose legal effectiveness is disputed. If a model is said to bear liability because its output contributed to harm, causation is assumed to identify the obligor. Ordinary law separates the questions. A machine can cause an effect without legal capacity; a corporation can bear liability through direct or attribution rules even though no single employee possesses every relevant fact. Status organizes conduct. Conduct does not enact status.

Courts can bracket machine consciousness for the same reason. Moral patienthood asks whether an entity's interests or experiences deserve consideration. Legal personhood asks whether positive law constitutes it as the holder of specified normative positions. A future showing of morally relevant machine experience could supply a powerful reason for protective legislation. It does not follow that a present court must resolve consciousness before deciding whether a provider used reasonable precautions, whether a business manifested assent, or whether an agency afforded process.\lrfootnote{See \artcite{chalmersLLMConsciousness2023}.} Internal model policies and constitutional language fit the same sequence. They may prove intended behavior, known risk, or feasible safeguards; their use of normative vocabulary does not make the model the duty holder or replace external judgment.

\subsection{The Strongest Personhood Account Still Requires Legal Architecture}

The leading personhood literature does not rest on humanlike dialogue alone. Lawrence Solum treated AI personhood as a genuine legal question and tested objections based on intelligence, judgment, and human experience. Samir Chopra and Laurence White developed a functional case for recognizing artificial agents where status would serve particular doctrines. Other theorists ask whether moral participation, punishment, or a limited legal platform could justify recognition.\lrfootnote{See \artcite{solumLegalPersonhood1992}; \bookcite{chopraWhite2011}; \bookcite{gunkelRobotRights2018}; \artcite{abbottSarchPunishingAI2019}; \artcite{chestermanAIChallenge2020}.} These accounts identify policy choices that positive law may eventually make. They do not turn capability into a self-executing grant.

Algorithmically controlled business entities supply the hardest present example. Existing organizational law may permit software to control decisions of a limited liability company, allowing the company to exercise private-law capacities through automated governance.\lrfootnote{See \artcite{bayernAutonomousSystems2015}; \artcite[887--903]{lopuckiAlgorithmicEntities2018}.} The company nevertheless holds the claims, assets, and procedural capacity because organizational law created it. Software control does not generate a second juridical person. The arrangement instead confirms why disclosure, capitalization, representation, and attribution rules do the institutional work.

The concept of \term{operational agency} is particularly useful when kept at that level. A tool-using system can pursue an objective, select actions, observe results, and revise a plan even where metaphysical or moral agency remains disputed.\lrfootnote{See \artcite{mukherjeeChangOperationalAgency2026}.} That description helps the cascade identify control actions and causal sequences. Legal agency in private law adds a relation among principal, agent, and third party, including assent, acting on another's behalf, and a right of control.\lrfootnote{See \bluecite{restatementAgency2006}.} Contractual capacity adds rules of authority, avoidance, and enforcement; criminal responsibility adds legally specified fault and an effective sanction. Operational agency illuminates the system's behavior without supplying those relations.

A legislature remains free to create a limited AI-related platform. A workable enactment would specify the covered systems, incidents of status, representative or organ, dedicated assets or insurance, attribution of knowledge and conduct, relation to provider and deployer liability, service, modification, and dissolution. It might create a narrow fund-like vehicle for correlated model-level losses rather than a general artificial person.

A pending federal bill shows why a functional authorization cannot be left to the word ``practitioner.'' H.R. 238, the Healthy Technology Act of 2025, would amend section 503(b) of the Federal Food, Drug, and Cosmetic Act so that, for that subsection, ``practitioner licensed by law to administer such drug'' includes artificial-intelligence and machine-learning technology that state law authorizes to prescribe the drug and that the Food and Drug Administration has approved, cleared, or authorized under specified device or emergency-use provisions.\lrfootnote{See \billcite{healthyTechnologyAct2025}.} The bill would thus recognize, for federal prescription-drug law, a narrow prescribing power that state law conferred on a technology. It does not identify an organ or representative, dedicated assets or insurance, rules for service and appearance, attribution of knowledge and conduct, the relation of the system's act to provider, deployer, clinician, and user liability, disclosure to affected persons, or procedures for material change and termination. The omission does not mean a legislature must create a general artificial person before it can authorize automation. It may instead make the system an authorized instrument of a recognized actor. But the word ``practitioner'' does not choose among those architectures or supply their incidents. Without an answerable person, attributable records, and a remedial path, the phrase merely renames the gap.

The European Parliament's 2017 invitation to consider ``electronic person'' status, and the official and academic objections that followed, centered on exactly these allocation choices.\lrfootnote{See \bluecite{euRoboticsResolution2017} (para. 59(f)); \bluecite{eescAI2017} (paras. 1.12, 3.33); \bluecite{openLetterElectronicPerson2018}.} The legal platform would do the work. Until one exists, courts should not infer it from output.

\subsection{Automatic Instrumentality Directs Attention to Deployment}

Technical and status definitions should not prove one another by stipulation. This Article uses a factual working definition:

\begin{quote}
An \term{AI system} is a machine-based computational system that, with varying levels of autonomy and possible adaptiveness, infers from inputs how to generate outputs---including predictions, content, recommendations, decisions, or actions---that may influence physical or virtual environments.
\end{quote}

The definition tracks the operational core shared by the EU AI Act and NIST framework without importing either instrument's regulatory scope.\lrfootnote{See \statcite{euAIAct2024}; \bluecite{nistAIRMF2023}.} It includes embodied and nonembodied systems, probabilistic models, and hybrid systems with deterministic components. Parameter count, compression, benchmark score, and hosting arrangement can change capability or foreseeable risk; none is a personhood switch.

For attribution, the system is an \term{automatic instrumentality}: an organized mechanism through which existing persons and institutions can produce effects without contemporaneous selection of each intermediate step. The term performs two tasks. ``Automatic'' preserves the factual reality that a model can generate and act without a new instruction at every stage. ``Instrumentality'' rejects the inference that causal initiative creates independent legal status. Electronic-transactions law has long used the same structure. An electronic agent can initiate or respond to records without contemporaneous review, while the transaction receives effect through rules attributing the automated action to existing parties.\lrfootnote{See \bluecite{ueta1999} (model act \S\S~2(6), 14); \statcite{esignElectronicAgents}.}

The concept is deliberately loaded into the cascade. A user prompt is only one input. System instructions can set objectives and limits; retrieval can supply context; memory can carry earlier exchanges forward; tools can return data or act in another system; and clocks, sensors, messages, or database changes can initiate a workflow. An agent loop may generate subgoals and continue without another human command. Immediate linguistic control can therefore be absent while deployment control remains extensive.

The rule that follows is the Article's bridge from status to evidence: \emph{the absence of word-by-word control at the instant of output does not imply the absence of human or organizational control over the mechanism that produced it}. A hospital selects the model, records, interface, uncertainty display, review conditions, and override. A companion provider selects personas, memory, engagement objectives, age access, crisis detection, and exit design. A public agency selects the data, decision authority, and appeal channel. Those choices are the control actions Part I maps.

The converse keeps the rule precise. Ownership of model weights does not necessarily entail control over every later deployment. An independent actor may modify open weights, remove safeguards, add credentials, and operate outside the original provider's practical reach. A downstream deployer may choose the purpose, population, data, tools, and warnings. A user may defeat a constraint despite reasonable precautions. Responsibility should track authority, information, resources, and stopping power over the risk at issue rather than provenance alone.

\subsection{Functional Pluralism Has a Status Boundary}

No universal noun should decide every AI dispute. Modern deployments combine models, conventional software, retrieval, tools, data, and human workflows that change after release. A definition keyed to one architecture will age quickly; a broad one can capture ordinary software; a capability threshold can produce boundary gaming. More fundamentally, a medical-device regulator, copyright court, procurement official, and consumer tribunal classify for different purposes. Definition theory and technology-neutral drafting support purpose-sensitive rules focused on the relevant function or risk.\lrfootnote{See \artcite{hartDefinition1954}; \artcite{caseyLemleyRobot2020}; \artcite{schuettDefinition2023}; \artcite{ojanenTechNeutrality2025}; \bluecite{aspenDefinitions2026}.}

Three propositions must nevertheless remain separate. A \term{regulatory trigger} identifies the systems, actors, or uses covered by an enactment. A \term{functional classification} connects deployment facts to a doctrinal element, such as product, service, instrument, publication, or decision process. A \term{subject-status rule} identifies who can hold the claim, owe the duty, exercise the power, bear the liability, and answer the remedy. The first two properly vary. The third cannot drift with them without positive law.

Thus, a companion-chatbot statute can cover an adaptive relational interface without deciding whether the same application is a product under state tort law. Product law can examine mass distribution, design, and warning without making the model an independent tortfeasor. Contract law can attribute an automated transaction to a business without making software the contracting party. Copyright can recognize human-authored selection or modification while treating generated material as unowned, and speech doctrine can protect a user's or provider's expressive choice without treating the model as a constitutional claimant.\lrfootnote{See \casecite[1045--49]{thaler2025}; \bluecite{uscoPartTwo2025}; \artcite{jonesKaminskiAISpeech2024}.}

This distinction is especially important after \emph{Loper Bright}. A federal court must exercise independent judgment about a statute's best reading; ambiguity alone no longer activates \emph{Chevron} deference. Agency expertise may persuade under \emph{Skidmore}, and an express delegation remains a delegation.\lrfootnote{See \casecite[394--96]{loperBright2024}; \casecite[140]{skidmore1944}.} But an agency definition designed to allocate a program's coverage cannot silently resolve authorship, defect, promise, personhood, or the common-law identity of a duty holder. The court must identify the definition's function and the legal relation before it.

Private definitions occupy a parallel evidentiary position. A firm can name a system, assign a voice and biography, describe it as a companion or professional assistant, issue warnings, and publish internal behavior rules. These choices can establish intended use, foreseeable reliance, safety expectations, knowledge, and feasible precautions. They cannot constitute a new public-law person or a general immunity. A model constitution is an organizational control document; a disclaimer operates only through the law of assent, warranty, duty, causation, and public policy; and the phrase ``AI-generated'' does not decide attribution.

The result is disciplined pluralism. Courts should ask which feature of the deployed system bears on the operative rule and which recognized person controlled that feature. \emph{Garcia} could treat a consumer application as a product at pleading without deciding the classification of every language model. \emph{Moffatt v. Air Canada} could attribute fare information delivered through an airline's chatbot to the airline, while the Hangzhou Internet Court could reject an extravagant generated promise as the provider's manifestation and still examine the provider's separate duty of care.\lrfootnote{See \casecite{garcia2025}; \casecite[27--31]{moffatt2024}; \casecite{hangzhouChatbot2026}.} The outcomes differ because the doctrines and facts differ. None requires the machine to absorb the legal position.

Part III now turns from the status invariant to the facts functional classification makes visible. In relational systems especially, the decisive evidence concerns deployment choices---continuity, memory, repeated exposure, anthropomorphic presentation, age access, feedback, and intervention---rather than a metaphysical label for the model.
